\section{Matched production and timing}
This production contains 8,000 bb0nu events at 2 mm pitch and 20 ms frames. The shared realization places ion drift uniformly between 5 and 145 cm, with complementary electron drift in the 150 cm chamber. Every event ID, primary-pair checksum and drift distance was checked against the existing matched common realization and 4 mm ion file. The electron readout was preserved. All events and signed response buckets are retained. The ion ledger is 897,306,359, exactly equal to the primary-pair total. The CSA used 2,878 nonlinear histories and 8 SPICE fallback histories among 2,232,910 pad histories. Production creates no SPICE control histories. Fresh filtered outputs were asserted before the run; actual circuit-deck, model, kernel-waveform and calibration checksums are saved in preflight.toml. Publication initially failed in the Julia sendfile copy. The oversized temporary was removed after verification of the corrected source. A bounded, checksum-verified copy completed publication with a 10% metadata-growth limit. Numerical timing covers the original calculation; publication recovery and registration are timed separately.
\begin{center}\small\begin{tabular}{p{11cm}r}\toprule Stage & Seconds\\\midrule
Ion diffusion and collection & 40.5933 \\
Covariance calibration setup & 64.7335 \\
Linear predictions and current preparation & 206.0798 \\
Nonlinear response and storage & 15.2782 \\
SPICE fallback workflow & 465.3983 \\
Three-frame reconstruction and summary & 376.7668 \\
Figures and presentation & 1.2545 \\
Event work excluding calibration and final presentation & 1063.5638 \\
Event work per event & 0.1329 \\
Publication recovery and registration & 200.2122 \\
\bottomrule\end{tabular}\end{center}
The numerical timer includes screening, nonlinear response, noise, reconstruction, output writing and consistency checks. Covariance calibration, final presentation and publication are separated. Event work includes fallback orchestration and its internal report generation when SPICE fallbacks occur. Stage and complete times are nested; the per-event value divides measured event work by 8,000.
